\documentclass[11pt]{article}
\usepackage[hmargin=.95in,top=.74in,bottom=.83in]{geometry}
\usepackage{amsmath,amssymb,amsthm}
\usepackage[hidelinks]{hyperref}
\newtheorem{theorem}{Theorem}[section]
\newtheorem{lemma}{Lemma}[section]
\newtheorem{proposition}{Proposition}[section]
\newcommand{\D}{\mathbb D}
\newcommand{\T}{\mathbb T}
\newcommand{\Z}{\mathbb Z}
\newcommand{\C}{\mathbb C}
\title{Independent domination roots of\\cozero-divisor graphs over residue rings}
\author{}
\date{Research note}
\hypersetup{pdftitle={Independent domination roots of cozero-divisor graphs over residue rings},
  pdfsubject={Research note on limiting roots of graph polynomials},pdfauthor={}}
\begin{document}
\maketitle
\begin{abstract}
\noindent
For a residue ring whose modulus is a power of one prime times a larger
prime, we describe the limiting nonzero roots of the independent
domination polynomial of its cozero-divisor graph as the exponent
increases. Except for the prime pair two and three, the root sets
converge in Hausdorff distance to the unit circle together with a
bounded discrete set outside it. These exterior points are
transcendental and accumulate at every point of the circle. They are
the reciprocals of the zeros in the open unit disk of an explicit gap
series. We characterize equality in an upper bound on their moduli
and obtain angular equidistribution for the nonzero roots counted with
multiplicity. For the exceptional pair the unique nonzero root equals
the negative of the exponent.

\end{abstract}
\begin{center}\small
\textit{Keywords.} Cozero-divisor graph, independent domination polynomial,
limiting zeros, lacunary power series, transcendence.\\
\textit{2020 Mathematics Subject Classification.} 05C31, 05C69, 30B10, 30C15.
\end{center}

\section{The limiting set}

For a finite commutative ring \(R\), the cozero-divisor graph
\(\Gamma'(R)\) has the nonzero nonunits as its vertices.
Two distinct vertices \(x,y\) are joined when \(x\notin Ry\) and
\(y\notin Rx\). An independent dominating set of a finite graph is
an inclusion-maximal independent vertex set. Its cardinality generates
one term in the independent domination polynomial
\(D_i(G,z)=\sum_{S} z^{|S|}\), where \(S\) ranges over these sets.

Fix primes \(p<q\), and write \(D_m(z)=D_i(\Gamma'(\Z/(p^m q)),z)\),
\(m\ge1\). Rather computed this polynomial in \cite[Theorem~2.4]{rather24}.
In \cite{rather26} he obtained an annular bound on its nonzero roots
for \((p,q)\ne(2,3)\) and asked for the set of their limit points as
\(m\to\infty\). A point is a limit point here if it is the limit of
roots chosen along some strictly increasing sequence of indices.
We determine this set and show that the finite root sets converge to
it in Hausdorff distance for each nonexceptional pair.

Write \(\D=\{w\in\C:|w|<1\}\), \(\T=\partial\D\), and let \(A_m\)
be the set of distinct nonzero roots of \(D_m\). For \(c=q-p-1>0\), put
\begin{equation}\label{series}
 e_j=1+cp^j,\qquad
 F(w)=1+\sum_{j\ge0}w^{e_j}\quad(w\in\D),\qquad
 E=\{w^{-1}:w\in\D,\ F(w)=0\}.
\end{equation}

\begin{theorem}\label{limit}
If \((p,q)=(2,3)\), then \(A_m=\{-m\}\) and there is no finite
limit point. For every other pair, \(A_m\) converges in Hausdorff
distance to the compact set \(\T\cup E\). The set \(E\) consists of
isolated transcendental numbers of modulus greater than one, and its
set of accumulation points is exactly \(\T\).
\end{theorem}

\noindent For example, when \((p,q)=(2,5)\) the unique limit of greatest
modulus is \(-1/r\), where \(r^3+r^5+r^9+r^{17}+\cdots=1\) and
\(0<r<1\). Thus an exterior limit can persist at positive distance
from the circle even though its share of the zeros tends to zero.
Proposition~\ref{radius} gives the corresponding radius bound for all
pairs.

\section{The reciprocal form of the graph polynomial}

In the generic \(i\)-term of the display in \cite[Theorem~2.4]{rather24},
the factor multiplying \(p^{m-1}-p^{i-1}\) should read \(q-1\),
as in the adjacent terms. We derive the reciprocal form directly
from the divisor chains.

\begin{lemma}\label{chains}
Put \(H=(q-1)p^{m-1}\) and \(c=q-p-1\). Then
\[
 D_m(z)=z^H F_m(1/z),\qquad F_m(w)=1+\sum_{j=0}^{m-1}w^{1+cp^j}.
\]
The degree of \(D_m\) is \(H\), and its zero at the origin has
multiplicity \(p^m-1\).
\end{lemma}
\begin{proof}
Let \(N=p^m q\). The vertices whose gcd with \(N\) equals \(d\)
form a class of size \(\varphi(N/d)\) for \(1<d<N\), \(d\mid N\).
Two principal ideals generated by vertices are comparable exactly when
their gcds with \(N\) are comparable under divisibility. A maximal
independent set therefore contains the full classes along one maximal
chain of proper nontrivial divisors. Conversely, every such chain
gives a maximal independent set.

Adjoining the endpoints \(1,N\), this chain takes \(m\) steps
multiplying by \(p\) and one step multiplying by \(q\). Let \(k\) be
the number of \(p\)-steps before the \(q\)-step. For \(k=m\) the
class sizes sum to \((q-1)\sum_{a=1}^{m}\varphi(p^{m-a})=H\). For
\(0\le k<m\), they sum to
\begin{align*}
(q-1)\sum_{a=1}^{k}\varphi(p^{m-a})
 + \sum_{a=k}^{m-1}\varphi(p^{m-a})
 &= (q-1)(p^{m-1}-p^{m-k-1})+p^{m-k}-1\\
 &=H-1-cp^{m-k-1}.
\end{align*}
Summing the resulting \(m+1\) monomials proves the identity, read
after expansion as a polynomial identity. Its smallest exponent is
\(H-1-cp^{m-1}=p^m-1\).
\end{proof}

The case \(c=0\) occurs only at \((p,q)=(2,3)\), when
\(D_m(z)=z^{2^m-1}(z+m)\). In particular \(m=4\) recovers
\cite[Example~2.2]{rather24}. Henceforth assume \(c>0\). Both
\(F_m\) and its monic reciprocal have constant coefficient one and all
other coefficients in \(\{0,1\}\). The elementary Cauchy bound applied
to each places \(A_m\) in the compact annulus \(1/2\le |z|\le2\).

\section{Zeros of the sections}

\begin{proof}[Proof of Theorem~\ref{limit}]
The series in \eqref{series} converges uniformly on compact subsets of
\(\D\). Any limit in \(\D\) of zeros of \(F_m\) must be a zero of \(F\).
Conversely, every disk isolating such a zero, with closure in \(\D\)
and no zeros on its boundary, contains exactly its multiplicity in
zeros of \(F_m\) for all sufficiently large \(m\), by Rouch\'e's
theorem. Inversion gives the description in \(|z|>1\).

Fix \(R>1\). For \(|w|\ge R\) and \(m\ge2\),
\[
 \left|\frac{F_m(w)}{w^{e_{m-1}}}-1\right|
 \le m R^{-(e_{m-1}-e_{m-2})}
 =mR^{-c(p-1)p^{m-2}}\longrightarrow0 .
\]
Thus \(A_m\) eventually avoids every closed disk \(|z|\le1/R\).

To describe the boundary of \(E\), we use Fuchs's sector theorem
\cite{fuchs}, stated in \cite[p.~515]{sons}. An analytic series on
\(\D\) with nonconstant exponents \(n_j\) satisfying
\(\inf_j n_{j+1}/n_j>1\) and corresponding coefficients \(a_j\)
satisfying \(\limsup_j|a_j|>0\) assumes each finite value infinitely
often in every open sector of \(\D\). Here \(a_j=1\) and
\(e_{j+1}/e_j\ge(1+cp)/(1+c)>1\). In each sector there are therefore
infinitely many zeros of \(F\), and only finitely many can belong to
any given compact subdisk. Taking sectors shrinking about a
prescribed angle shows that every point of \(\T\) is approached by
these zeros. Since \(F(0)=1\) and the zeros are isolated in \(\D\),
inversion makes \(E\) bounded and discrete in \(|z|>1\), with
accumulation set exactly \(\T\).

Let \(K=\T\cup E\). Compactness of the common annulus and the
classification of subsequential limits give
\[\sup_{a\in A_m}\operatorname{dist}(a,K)\longrightarrow0.\]
Every neighborhood
of a point of \(E\) meets \(A_m\) eventually, by Rouch\'e's theorem.
Every neighborhood of a point of \(\T\) contains a point of \(E\)
and hence also meets \(A_m\) eventually. For any \(\varepsilon>0\),
covering \(K\) by finitely many balls of radius \(\varepsilon/2\)
centered on \(K\) gives
\(\sup_{b\in K}\operatorname{dist}(b,A_m)\le\varepsilon\) eventually.
Together these statements prove Hausdorff convergence.

Finally write \(G_p(s)=\sum_{j\ge0}s^{p^j}\), so that
\(F(w)=1+wG_p(w^c)\). Mahler's theorem on the Fredholm series
\cite{mahler}, in the form stated in \cite[p.~269]{becker}, asserts
that \(G_d(s)\) is transcendental for every integer \(d\ge2\) and
every algebraic \(s\) with \(0<|s|<1\). If a zero \(w\in\D\) were
algebraic, then \(w^c\) would satisfy these hypotheses, while
\(G_p(w^c)=-1/w\) would be algebraic. Thus \(w\) and its reciprocal
are transcendental.
\end{proof}

\begin{proposition}\label{distribution}
For \(c>0\), the empirical measures of the nonzero roots of \(D_m\),
counted with multiplicity, converge weakly to normalized arclength
\(\sigma\) on \(\T\). For every \(p<q\), the empirical measures of
all roots converge weakly to
\[
 \frac{p}{q-1}\delta_0+\frac{q-p-1}{q-1}\sigma .
\]
\end{proposition}
\begin{proof}
For \(c>0\), \(F_m\) has degree \(d_m=1+cp^{m-1}\), monic leading
term, constant coefficient one and coefficient sum \(m+1\). Thus
\(\log(m+1)=o(d_m)\). The deterministic zero-distribution criterion
of Hughes and Nikeghbali \cite{hughes} gives radial concentration
near \(\T\) and angular equidistribution, hence weak convergence to
\(\sigma\). Inversion on the common annulus preserves the limit.
The fraction of all roots at zero is \((p^m-1)/H\to p/(q-1)\).
For \(c=0\) the single moving root has mass \(1/H\to0\), which gives
the stated formula also in this case.
\end{proof}

\section{The outer radius}

\begin{proposition}\label{radius}
Assume \(c>0\), and let \(r\in(0,1)\) be the unique solution of
\(\sum_{j\ge0}r^{e_j}=1\). Set
\[
 \Delta=c(p-1),\qquad g=\gcd(1+c,\Delta)=\gcd(1+c,p-1).
\]
Every limit point has modulus at most \(1/r\). The bound is attained
if and only if \(\Delta/g\) is even. When attained, it is attained at
exactly \(g\) points \(1/(ru)\), characterized by
\[
 u^\Delta=1,\qquad u^{1+c}=-1.
\]
Each corresponding zero \(ru\) of \(F\) is simple.
\end{proposition}
\begin{proof}
The defining sum increases continuously from zero to infinity as
\(r\) increases from zero to one. For a zero \(w\in\D\) of \(F\),
\(1\le\sum_{j\ge0}|w|^{e_j}\), so \(|w|\ge r\).
Equality holds precisely when, for \(u=w/r\), all \(u^{e_j}=-1\):
indeed \(\sum r^{e_j}=1\), and the real part of each \(u^{e_j}\)
is at least \(-1\). The differences \(e_j-e_0=c(p^j-1)\) are all
multiples of \(\Delta\), with \(e_1-e_0=\Delta\). Thus the displayed
pair of phase equations is necessary and sufficient, and conversely
any such \(u\) gives \(F(ru)=0\). The power map \(u\mapsto u^{1+c}\)
on the cyclic group of \(\Delta\)-th roots of unity has kernel
of order \(g\) and image of order \(\Delta/g\). Its image contains
\(-1\) precisely when that order is even. At a solution,
\[
 F'(ru)=-u^{-1}\sum_{j\ge0}e_jr^{e_j-1}\ne0.
\]
As \(r<1\), the unit circle supplies no further equality points.
\end{proof}

For \(p=2\) we have \(c\) even and \(g=1\), so \(u=-1\) gives the
unique outermost point. For odd \(p\) all exponents are even,
and no exterior limit is real. The radius criterion can then go either way:
for \((p,q)=(5,7)\) the phase equations give \(u=\pm i\),
while for \((p,q)=(3,7)\) the compact limiting set lies strictly
inside \(|z|=\smash{1/r}\).


\section*{Verification and disclosures}
The Lean~4 supplement checks the reciprocal-series estimates over
the complex numbers, local uniform convergence of the sections, the
exclusion of interior limiting roots, and the compactness argument
for Hausdorff convergence. A maximum-modulus argument in Lean gives
a zero of every sufficiently late section near each zero of the
limiting series. The written Rouch\'e argument additionally gives the
exact local multiplicity. Fuchs's sector value theorem and Mahler's
transcendental-value theorem enter as explicit hypotheses in the
formal deductions. The weak convergence of the empirical measures
uses the published criterion of Hughes and Nikeghbali in the written
proof. For the outer-radius criterion, the Lean certificate covers
the common divisors of the exponent differences.

OpenAI Codex was used in selecting the problem, developing the
mathematical arguments, searching the literature, drafting and
revising the text, and writing and testing the Lean and Python code.

The text of this research note is licensed under
\href{https://creativecommons.org/licenses/by/4.0/}{Creative Commons Attribution 4.0 International (CC BY 4.0)}.
For attribution, cite the title of this note and its source repository,
\url{https://github.com/xieboyan1/cozero-divisor-roots}.

\small
\begin{thebibliography}{00}
\bibitem{becker}
P.-G. Becker, Transcendence measures for the values of generalized
Mahler functions in arbitrary characteristic,
\emph{Publ. Math. Debrecen} \textbf{45} (1994) 269--282.
\href{https://doi.org/10.5486/PMD.1994.1437}{doi:10.5486/PMD.1994.1437}.
\bibitem{fuchs}
W.~H.~J. Fuchs, Topics in Nevanlinna theory, in \emph{Proceedings of
the NRL Conference on Classical Function Theory}, Naval Research
Laboratory, Washington, DC, 1970, 1--32.
\bibitem{hughes}
C.~P. Hughes and A. Nikeghbali, The zeros of random polynomials cluster
uniformly near the unit circle, \emph{Compos. Math.} \textbf{144}
(2008) 734--746. \href{https://doi.org/10.1112/S0010437X07003302}{doi:10.1112/S0010437X07003302}.
\bibitem{mahler}
K. Mahler, Arithmetische Eigenschaften der L\"osungen einer Klasse
von Funktionalgleichungen, \emph{Math. Ann.} \textbf{101} (1929)
342--366. \href{https://doi.org/10.1007/BF01454845}{doi:10.1007/BF01454845}.
\bibitem{rather24}
B.~A. Rather, Independent domination polynomial for the cozero divisor
graph of the ring of integers modulo \(n\), \emph{Discrete Math. Lett.}
\textbf{13} (2024) 36--43.
\href{https://www.dmlett.com/archive/v13/DML24_v13_pp36-43.pdf}{doi:10.47443/dml.2023.215}.
\bibitem{rather26}
B.~A. Rather, Zeros of the independent domination polynomial of the
cozero-divisor graphs of the ring \(\Z_n\), \emph{Siberian Math. J.}
\textbf{67} (2026) 1434--1443.
\href{https://doi.org/10.1134/S0037446626060194}{doi:10.1134/S0037446626060194}.
\bibitem{sons}
L.~R. Sons, Zeros of sums of series with Hadamard gaps,
\emph{Pacific J. Math.} \textbf{43} (1972) 515--524.
\href{https://msp.org/pjm/1972/43-2/pjm-v43-n2-p24-p.pdf}{Publisher's PDF}.
\end{thebibliography}
\end{document}
